\subsection*{問題}

\begin{enumerate}
  \item $\bigl(\mathbb{Z}[\omega]/(3)\bigr)^{\times}$ の位数を求めよ。
  \item $\bigl(\mathbb{Z}[\omega]/(9)\bigr)^{\times}$ を有限Abel群の直和で表せ。
\end{enumerate}

\subsection*{答案}

\begin{enumerate}
  \item
  $\mathbb{Z}[\omega]\simeq\mathbb{Z}[x]/(x^2+x+1)$ より、
  \begin{align*}
    \mathbb{Z}[\omega]/(3)
      &\simeq \mathbb{F}_3[x]/(x^2+x+1) \\
      &= \mathbb{F}_3[x]/(x-1)^2 \\
      &\simeq \mathbb{F}_3[t]/(t^2).
  \end{align*}
  これは局所環なので、極大イデアル $(t)/(t^2)$ の外の元がすべて単元。
  よって6個。

  \item
  \begin{align*}
    \mathbb{Z}/9\mathbb{Z}[x]/(x^2+x+1)
      &= \mathbb{Z}/9\mathbb{Z}[x]/\bigl((x-4)^2+3\bigr) \\
      &\simeq \mathbb{Z}/9\mathbb{Z}[t]=A.
  \end{align*}
  （$t^2=-3$）の位数は81。

  $\operatorname{nil}A\ni t,3,-3$ が分かるので、$\operatorname{nil}A$ がイデアルだから、
  $a\in\mathbb{Z}/9\mathbb{Z}$、$b\in3\mathbb{Z}/9\mathbb{Z}$ に対して
  \[
    at+b\in\operatorname{nil}A
  \]
  が分かる。
  逆に $at+b\in\operatorname{nil}A$ なら $b\in\operatorname{nil}A$ なので、
  $b\in3\mathbb{Z}/9\mathbb{Z}$。

  よって
  \[
    \#(\operatorname{nil}A)=27
  \]
  で、これらは単元ではない。

  逆に $b\notin3\mathbb{Z}/9\mathbb{Z}$ で $at+b$ を考える。
  \[
    (at+b)(at-b)=a^2t^2-b^2=-3a^2-b^2
      \notin3\mathbb{Z}/9\mathbb{Z}
  \]
  は単元なので、
  \[
    (at+b)(at-b)(-3a^2-b^2)^{-1}=1.
  \]
  よって $\bigl(\mathbb{Z}[\omega]/(9)\bigr)^{\times}$ は
  $81-27=54$ 個の元を持つAbel群なので、構造定理より
  \[
    \mathbb{Z}/54\mathbb{Z},\qquad
    \mathbb{Z}/3\mathbb{Z}\times\mathbb{Z}/18\mathbb{Z},\qquad
    \mathbb{Z}/3\mathbb{Z}\times\mathbb{Z}/3\mathbb{Z}\times\mathbb{Z}/6\mathbb{Z}
  \]
  のどれかになる。

  \begin{align*}
    (at+b)^3
      &=a^3(-3)t+3atb^2+b^3 \\
      &=3(a-a^3)t+b^3.
  \end{align*}
  ここで
  \[
    a-a^3=-a(a-1)(a+1)
  \]
  は（連続する3整数の積ゆえ）必ず $3\mathbb{Z}/9\mathbb{Z}$ に属すので、
  \[
    (at+b)^3=b^3.
  \]
  よって
  \[
    (at+b)^6=b^6=1
    \qquad\bigl(b\notin3\mathbb{Z}/9\mathbb{Z}\bigr).
  \]
  よって位数は高々6なので、
  \[
    \mathbb{Z}/3\mathbb{Z}\times\mathbb{Z}/3\mathbb{Z}\times\mathbb{Z}/6\mathbb{Z}
  \]
  が答え。
\end{enumerate}
